% --- Paquetes Necesarios ---

\usepackage[utf8]{inputenc}
\usepackage{titlesec}
\usepackage{changepage}
\usepackage{ragged2e} % Si falla, podemos comentar esta línea
\usepackage{caption}

\setlength{\parindent}{0pt}
\setlength{\parindent}{0pt}

% --- CONFIGURACIÓN DE TÍTULOS ---

% Título I: CENTRADO CON SALTO
\titleformat{\section}[display]
  {\normalfont\fontsize{12}{14}\bfseries\centering} % Quitamos \uppercase de aquí
  {CAPÍTULO \thesection}
  {1ex}
  {\MakeUppercase} % Lo aplicamos directamente al cuerpo del título
\titlespacing{\section}{0pt}{1ex}{2ex}

% Título II: IZQUIERDA
\titleformat{\subsection}
  {\normalfont\fontsize{12}{14}\bfseries}
  {\thesubsection}{1em}{}
\titlespacing{\subsection}{0pt}{1.5ex}{1ex}

% Título III: INDENTADO
\titleformat{\subsubsection}
  {\normalfont\fontsize{12}{14}\bfseries}
  {\thesubsubsection}{1em}{}
\titlespacing{\subsubsection}{1.5em}{1.5ex}{1ex}

% --- CONFIGURACIÓN DE PROFUNDIDAD DE NUMERACIÓN ---
\setcounter{secnumdepth}{4} % Permite numerar hasta el nivel 4 (tituloIV)
\setcounter{tocdepth}{4}    % Opcional: Incluye el nivel 4 en el índice (si lo deseas)
% Título IV: MÁS INDENTADO
\titleformat{\paragraph}
    {\normalfont\fontsize{12}{14}\bfseries}
    {\theparagraph} % <-- Esto es lo que imprime el 1.1.1.1
    {1em}
    {}
\titlespacing{\paragraph}{4.0em}{1.5ex}{1ex}

% --- COMANDOS PERSONALIZADOS ---
\newcommand{\tituloI}[1]{\section{#1}}
\newcommand{\tituloII}[1]{\subsection{#1}}
\newcommand{\tituloIII}[1]{\subsubsection{#1}}
\newcommand{\tituloIV}[1]{\paragraph{#1}} 
% El \mbox{}\\} es para asegurar que el párrafo empiece en la línea siguiente si así lo deseas

% --- ENTORNO PARA PÁRRAFOS (Versión Estable) ---
\newenvironment{parrafo}[1]{
    \begin{adjustwidth}{#1}{0pt}
    %\justifying % Si da error de "Undefined control sequence", borrar esta línea
    
    \setlength{\parindent}{0pt}
}{
    \end{adjustwidth}
}

% --- COMANDOS PARA PÁRRAFOS SIN SANGRÍA (CORREGIDOS) ---

% Párrafo para Título I
\newcommand{\parrafoI}[1]{%
    \noindent\ignorespaces#1\par\vspace{1em}%
}

% Párrafo para Título II (Margen 2.5em)
\newcommand{\parrafoII}[1]{%
    \begin{adjustwidth}{2.5em}{0pt}%
        \setlength{\parindent}{0pt}%
        \noindent\ignorespaces#1\par%
    \end{adjustwidth}%
}

% Párrafo para Título III (Margen 5em)
\newcommand{\parrafoIII}[1]{%
    \begin{adjustwidth}{5em}{0pt}%
        \setlength{\parindent}{0pt}%
        \noindent\ignorespaces#1\par%
    \end{adjustwidth}%
}
% Párrafo para Título IV (Margen 6.5em o 7.0em según tu preferencia de escalado)
\newcommand{\parrafoIV}[1]{%
    \begin{adjustwidth}{8.4em}{0pt}%
        \setlength{\parindent}{0pt}%
        \noindent\ignorespaces#1\par%
    \end{adjustwidth}%
}
% --- LISTAS PERSONALIZADAS POR NIVEL ---

% Lista para Nivel I (Margen 0)
\newlist{listaI}{itemize}{1}
\setlist[listaI]{
    label=\textbullet,
    nosep,
    leftmargin=1.5em, % Espacio para que la viñeta no se salga al margen
    topsep=0.5em,     % Espacio de 0.5em con el título
    after=\vspace{1em}
}

% Lista para Nivel II (Alineada con Parrafo II)
\newlist{listaII}{itemize}{1}
\setlist[listaII]{
    label=\textbullet,
    nosep,
    leftmargin=4.0em, % 2.5em del parrafo + 1.5em para la viñeta
    topsep=0.5em,
    after=\vspace{1em}
}

% Lista para Nivel III (Alineada con Parrafo III)
\newlist{listaIII}{itemize}{1}
\setlist[listaIII]{
    label=\textbullet,
    nosep,
    leftmargin=5.8em, % 4.5em del parrafo + 1.5em para la viñeta
    topsep=0.5em,     % Aquí controlamos la distancia exacta al título
    after=\vspace{1em}
}

% Definición corregida para listas numeradas de Nivel III
\newlist{listaIIInum}{enumerate}{1} % Cambiado itemize por enumerate
\setlist[listaIIInum]{
    label=\arabic*.,      % El asterisco indica que use el contador de esta lista
    nosep,
    leftmargin=6.2em,     
    topsep=0.5em,         
    after=\vspace{1em}
}

% Lista para Nivel III (Alineada con Parrafo III)
\newlist{listaIV}{itemize}{1}
\setlist[listaIV]{
    label=\textbullet,
    nosep,
    leftmargin=9.5em, % 4.5em del parrafo + 1.5em para la viñeta
    topsep=0.5em,     % Aquí controlamos la distancia exacta al título
    after=\vspace{1em}
}

\newlist{listaOE}{enumerate}{1}
\setlist[listaOE]{
    label=\textbf{OE\arabic*}, % Define el formato OE1, OE2...
    leftmargin=7.7em,             % Alineado con el margen del Título III
    labelsep=1em,               % Espacio entre el código y el texto
    nosep                       % Elimina espacios verticales excesivos
}

\newlist{listaPE}{enumerate}{1}
\setlist[listaPE]{
    label=\textbf{PE\arabic*}, % Define el formato PE1, PE2...
    leftmargin=7.7em,             % Alineado con el margen del Título III
    labelsep=1em,               % Espacio entre el código y el texto
    nosep                       % Elimina espacios verticales excesivos
}

\newlist{listaVar}{enumerate}{1}
\setlist[listaVar]{
    label=\textbf{VAR\arabic*}, % Define el formato PE1, PE2...
    leftmargin=7.7em,             % Alineado con el margen del Título III
    labelsep=1em,               % Espacio entre el código y el texto
    nosep                       % Elimina espacios verticales excesivos
}

\newlist{listaHE}{enumerate}{1}
\setlist[listaHE]{
    label=\textbf{HE\arabic*}, % Define el formato PE1, PE2...
    leftmargin=8em,             % Alineado con el margen del Título III
    labelsep=1em,               % Espacio entre el código y el texto
    nosep                       % Elimina espacios verticales excesivos
}

\newlist{listaVD}{enumerate}{1}
\setlist[listaVD]{
    label=\textbf{VD\arabic*}, % Define el formato PE1, PE2...
    leftmargin=8em,             % Alineado con el margen del Título III
    labelsep=1em,               % Espacio entre el código y el texto
    nosep                       % Elimina espacios verticales excesivos
}

% Tablas
\captionsetup[table]{
    name=Tabla,
    labelsep=period,
    justification=centering,
    singlelinecheck=false
}

% Figuras
\captionsetup[figure]{
    name=Figura,
    labelsep=period,
    justification=centering,
    singlelinecheck=false
}